The measured width contrast rejects the strong claim that width mainly accelerates memorization while leaving the slow process unchanged. The change occurs after both widths can fit the same training size. It is compatible with a feature-discovery bottleneck: added width changes optimization access to useful structured solutions under regularization. However, no internal circuit was isolated, and alternatives involving conditioning, redundancy, and Adam's effective updates remain possible.

Persistence of the contrast under expected initial norm matching argues against initial total norm alone as a sufficient explanation. It does not refute norm-based theories that depend on the entire optimization path or circuit efficiency.

The direction intervention supplies additional evidence: fixing hidden directions at step 1,000 preserves final training interpolation yet removes all five reference generalization events. Readout learning, individual feature magnitudes, and weight decay remain active, so their constrained continuation is insufficient from that state under these updates. Fixing directions only at step 10,000 retains 3/5 events, showing that continued direction changes are not required after that point in every observed run. The intervention still changes the optimization path, and identifies neither a unique circuit nor a general necessity theorem.

The operational distinction between occurrence and latency matters. A configuration with no confirmed event cannot be assigned a small delay; a configuration with both rapid fitting and later generalization can exhibit a large ratio even when its absolute waiting time decreases. The success map, restricted waiting time, and delayed-event counts therefore answer complementary questions. Finite-horizon success is an experimentally useful setting-selection criterion, but it is not a universal phase boundary.
